% =====================================================================
%  2bat-ud6-2.tex  —  SESSIÓ 2: Què ens diu una mostra?
%  (sense preàmbul: s'inclou des de main.tex)
% =====================================================================
\horatitol{Sessió 2 — Què ens diu una mostra?}%
{Activitat 1: valorar si la mostra d'un estudi real és representativa de la població i detectar-ne possibles biaixos de selecció.}

\subsection{Idea clau}

Treballem en parelles el cas de l'Activitat 1: una enquesta sobre un col·lectiu, amb la
informació de \textbf{com s'ha triat la mostra}. La feina no és calcular res, sinó
\textbf{valorar} si aquesta mostra és representativa de tota la població que es vol
estudiar, o si deixa fora algun grup important.

\begin{idea}
\textbf{Representativitat i biaix.}
\begin{itemize}[leftmargin=1.4em, itemsep=2pt]
  \item Una mostra és \textbf{representativa} si reflecteix la diversitat real de la
        població: els grups que hi ha al tot hi són en una proporció semblant.
  \item Hi ha \textbf{biaix de selecció} quan la manera de triar la mostra deixa fora (o
        infrarepresenta) algun grup, de manera sistemàtica.
  \item \textbf{Condició clau:} la inferència només és vàlida si la mostra és
        representativa. Una conclusió és tan bona com la mostra de la qual prové.
\end{itemize}
\textbf{La pregunta a fer-se sempre:} qui queda fora d'aquesta mostra? Si en queda fora un
grup rellevant, el resultat pot estar esbiaixat.
\end{idea}

\subsection{Exemples}

\begin{exemple}[una enquesta només per internet]
Un estudi sobre les necessitats de les persones grans d'un barri es fa amb una enquesta
\textbf{en línia}. Sembla còmode, però deixa fora, sistemàticament, qui no té accés digital
—sovint les persones més grans i vulnerables, que són justament el col·lectiu d'interès.
La mostra \textbf{no és representativa}: hi ha un \textbf{biaix de cobertura}, i el resultat
descriuria només les persones grans connectades, no totes.
\end{exemple}

\begin{exemple}[preguntar només a qui continua venint]
Per avaluar la satisfacció amb un servei, s'enquesta només qui hi \textbf{continua anant}.
Qui l'ha abandonat —potser per insatisfacció— no respon. La mostra sobreestima la
satisfacció: és un \textbf{biaix de selecció} que empeny el resultat cap amunt respecte del
valor real.
\end{exemple}

\subsection{Exercicis resolts}

\begin{exresolt}[detectar qui queda fora]
Una enquesta sobre l'ús del transport públic es fa \textbf{a les parades d'autobús} en
hora punta. És representativa de tota la població? Qui en queda fora?
\solucio
\textbf{No} és representativa de tota la població: només arriba a qui \emph{ja} fa servir
el bus (i en hora punta). En queden fora els qui no l'utilitzen (perquè van a peu, en
cotxe, o no es mouen), i també qui viatja en altres horaris. El resultat sobreestimaria
l'ús del transport públic: hi ha biaix de selecció.
\resposta{No; en queda fora qui no fa servir el bus. Sobreestimaria l'ús.}
\end{exresolt}

\begin{exresolt}[com millorar la mostra]
En l'estudi de les persones grans del barri (enquesta en línia), com es podria fer la
mostra més representativa?
\solucio
Cal arribar també a qui no té accés digital: per exemple, combinant l'enquesta en línia
amb entrevistes \textbf{presencials} (a casals, centres de dia, a domicili) i telefòniques.
Així la mostra inclouria tant les persones connectades com les que no ho estan, i
reflectiria millor la diversitat real del col·lectiu.
\resposta{Afegint vies presencials i telefòniques per incloure qui no té accés digital.}
\end{exresolt}

\subsection{Exercicis per fer}

\begin{exercicis}
\begin{exlist}
  \item Una enquesta sobre hàbits de lectura es fa \textbf{a la sortida d'una biblioteca}.
        És representativa de tota la població? Qui en queda fora?
  \item Per conèixer l'opinió dels veïns d'un barri, es reparteix un formulari
        \textbf{només en català}. Quin grup podria quedar infrarepresentat?
  \item Un estudi sobre l'ús d'un servei juvenil enquesta els joves \textbf{un dissabte al
        matin} en un skatepark. Digues un biaix possible i en quina direcció esbiaixaria
        el resultat.
  \item Explica amb les teves paraules per què una conclusió «és tan bona com la mostra de
        la qual prové».
  \item \textbf{(Millora)} Proposa com faries més representativa la mostra de l'exercici 1
        (enquesta a la sortida de la biblioteca).
  \item \textbf{(Interpretació)} Un estudi conclou que «la majoria de la població fa
        esport regularment», però la mostra es va recollir \textbf{en un gimnàs}. Fins on
        es pot generalitzar aquesta conclusió? Amb quina cautela?
\end{exlist}
\end{exercicis}

% ---------------------------------------------------------------------
\fullsolucions{Sessió 2}
\begin{solucions}
\begin{sollist}
  \item No: només arriba a qui va a la biblioteca (persones ja lectores). En queda fora
        qui no hi va. Sobreestimaria els hàbits de lectura de la població.
  \item Podria quedar infrarepresentada la població que no llegeix català amb fluïdesa
        (per exemple, persones nouvingudes), de manera que la seva opinió quedaria fora.
  \item Biaix de selecció: s'arriba només a un perfil concret de joves (els que van a
        l'skatepark el dissabte). Esbiaixaria el resultat cap als interessos i l'ús
        d'aquest grup, no cap a tots els joves.
  \item Resposta oberta. Idea: si la mostra deixa fora grups o està mal triada, el
        resultat que se n'obté no reflecteix la població; per bé que es calculi, la
        conclusió arrossega el biaix de la mostra.
  \item Resposta oberta. Idea: combinar diferents llocs i moments (no només la biblioteca)
        i incloure persones que no hi van, per exemple amb una enquesta a domicili o en
        diversos punts del barri triats a l'atzar.
  \item Resposta oberta. Idea: no es pot generalitzar a tota la població, perquè en un
        gimnàs gairebé tothom fa esport (mostra esbiaixada); la conclusió només valdria,
        com a molt, per als usuaris de gimnasos, i cal dir-ho explícitament.
\end{sollist}
\end{solucions}
